%\myChapterStar{Titre}{Titre court}{Ajouter à la table des matières? (false|true|chapter|section|subsection|subsubsection -chapter par défaut-)}
\myChapterStar{Conclusion générale}{}{true}
\myMiniToc{}{Sommaire}

Ce mémoire présente une recherche innovante à l'intersection de deux domaines informatiques distincts: les Grammaires Attribuées Gardées (GAG) comme paradigme de composition incrémentale des services, et la transformée de Fourier rapide (FFT) comme cas d'usage benchmark pour l'évaluation des performances.

\mySectionStar{Problématique étudiée et choix méthodologiques}{}{true}

La question de recherche principale porte sur l'évaluation de l'efficacité de la composition incrémentale dans un contexte de calcul intensif, dans quelle mesure les services paresseux peuvent-ils constituer une alternative viable aux approches traditionnelles pour l'implémentation d'algorithmes computationnellement exigeants, tout en préservant les avantages de la composition incrémentale? À cet effet, nous avons d'abord modélisé un algorithme de FFT paresseux au moyen de GAG, permettant ainsi de dissocier la spécification déclarative (règles de productions) de l'exécution (fonctions locales). Ensuite, ces services ont été déployés dans un environnement conteneurisé (Docker) orchestré par Kubernetes sur un cluster. Enfin, nous avons élaboré un protocole expérimental mesurant la latence en fonction de la taille des transformées afin d’analyser l’impact de l’incrémental sur les performances.

\vspace{1cm}
\mySectionStar{Analyse critique des résultats obtenus}{}{true}

Les résultats expérimentaux révèlent que, d'une part, la latence reste maîtrisée pour des FFT de grande taille du fait que l'incrémentation ne concerne que les portions modifiées, et d'autre part, la scalabilité s'améliore grâce à la distribution fine des tâches FFT sur plusieurs nœuds Kubernetes sans recomposition globale. Par ailleurs, l'approche paresseuse limite l'utilisation mémoire en évitant un recalcule exhaustif, ce qui permet de rivaliser avec des compositions statiques optimisées. Toutefois, il convient de noter que pour des transformées de très petite taille, l'overhead initial du moteur GAG peut surpasser les gains apportés, et que la performance dépend étroitement de la granularité choisie dans les GAG, rendant critique un paramétrage adapté.

\mySectionStar{Quelques perspectives}{}{true}

Ce travail propose une approche certes prometteuse mais pas parfaite, c'est pourquoi plusieurs perspectives de recherche se posent: 
% Ce travail propose une approche certes prometteuse, mais encore perfectible, ce qui ouvre plusieurs perspectives de recherche :
\begin{itemize}
    \item Étendre l’expérimentation à des infrastructures haut de gamme (bare‑metal ou instances cloud) pour valider la scalabilité et la robustesse en conditions réelles;
    \item Assurer la récupération automatique des données perdues en cas de panne de service;
    \item Exploiter la parallélisation de l’opération de combinaison en l’exprimant sous forme matricielle pour mettre en lumière la pertinence des services paresseux sur des modèles de calcul parallèle;
    \item Développer un module de supervision adaptative capable d’ajuster dynamiquement la granularité d’incrémentation en fonction de la charge des nœuds, afin de maximiser l’utilisation des ressources et d’optimiser le coût opérationnel;
    \item Étendre l’usage de la composition incrémentale à des algorithmes non nécessairement de type DPR, mais qui présentent des structures de dépendances explicites entre leurs composants.% L’objectif serait de construire un graphe de calcul reflétant ces relations, permettant ainsi d’activer uniquement les services réellement affectés par une modification, et de paralléliser leur exécution de manière ciblée.
\end{itemize}

\myCleanStarChapterEnd

\myCleanStarChapterEnd
